\begin{figure}[t]
\centering
\colorlet{cadv}{red!80!black}
\colorlet{chon}{blue!70!black}
\colorlet{call}{violet}
% rho = 1/2. Rays of slope 3/4 (everybody), 1/2 (honest) and 1/4 (adversary),
% each inside a band of slopes (1 +- 0.12) times its own. The client receives
% C of height 3.2 at round r = 8. The worst case forks at round 2.4, height
% 1.8: 0.75*2.4 = 1.8 and 1.8 + 0.25*(8-2.4) = 3.2. Our rule keeps height
% 1.8, the simple rule keeps 3.2 - 0.25*8 = 1.2.
\begin{tikzpicture}[x=0.83cm,y=0.7cm,>=stealth,font=\footnotesize]
  \def\xe{8.8}
  % Chernoff bands
  \fill[call,opacity=0.13] (0,0) -- (\xe,0.84*\xe) -- (\xe,0.66*\xe) -- cycle;
  \fill[chon,opacity=0.13] (0,0) -- (\xe,0.56*\xe) -- (\xe,0.44*\xe) -- cycle;
  \fill[cadv,opacity=0.13] (0,0) -- (\xe,0.28*\xe) -- (\xe,0.22*\xe) -- cycle;
  % axes
  \draw[->] (0,0) -- (10.8,0) node[below left] {round};
  \draw[->] (0,0) -- (0,7) node[left] {height};
  % rates
  \draw[thick,call] (0,0) -- (\xe,0.75*\xe)
    node[pos=0.72,above,sloped] {everybody};
  \draw[thick,chon] (0,0) -- (\xe,0.5*\xe)
    node[pos=0.72,above,sloped] {honest};
  \draw[thick,cadv] (0,0) -- (\xe,0.25*\xe)
    node[pos=0.45,below,sloped] {adversary};
  % guides
  \draw[lim] (8,0) -- (8,4);
  \draw[lim] (0,3.2) -- (8,3.2);
  \draw[lim] (0,1.8) -- (8.95,1.8);
  \draw[lim] (0,1.2) -- (10.75,1.2);
  \draw[lim] (8,3.2) -- (10.75,3.2);
  % the received chain: its prefix C_0 along the top rate, then the private tail
  \draw[very thick,call] (0,0) -- (2.4,1.8);
  \draw[very thick,dashed,cadv] (2.4,1.8) -- (8,3.2);
  \fill (2.4,1.8) circle (1.8pt) node[above left] {$\ch_0$};
  \fill (8,3.2) circle (1.8pt) node[above left] {$\ch$};
  \fill (8,4) circle (1.8pt) node[above left] {$\ch'$};
  \draw[very thick] (7.92,1.8) -- (8.08,1.8)
    node[below right=-1pt] {$\prune{\ch}{\prn(r,\hgt)}$};
  \draw[very thick] (7.92,1.2) -- (8.08,1.2)
    node[below right=-1pt] {$\prune{\ch}{\prn(r)}$};
  % what the two rules prune
  \draw[brc] (9,1.8) -- (9,3.2)
    node[midway,right=5pt,align=left] {our rule\\$\prn(r,\hgt)$};
  \draw[brc] (10.8,1.2) -- (10.8,3.2)
    node[midway,right=5pt,align=left] {simple rule\\$\prn(r)$};
  % what our rule keeps
  \draw[brc] (-0.1,1.8) -- (-0.1,0) node[midway,left=5pt,rotate=90,anchor=south]
    {safe prefix};
  \node[left] at (0,3.2) {$\hgt$};
  \node[below] at (8,0) {$r$};
\end{tikzpicture}
\caption{The pruning depth, drawn for $\advr=\frac{1}{2}$ as height against
round. The three lines are the mining rates of everybody together (violet), of
the honest validators (blue) and of the adversary alone (red). The band
around each line is how far a typical execution can stray from it. At round
$r$ the client receives the chain $\ch$ of height $\hgt$, while an honest
validator holds the longer chain $\ch'$. In the worst case, $\ch$ grows at the
top rate up to its prefix $\ch_0$, and the adversary mines the rest in
private, along the dashed red line. Our rule prunes $\prn(r,\hgt)$ blocks,
which takes $\ch$ down to the height of $\ch_0$. The simple rule prunes
$\prn(r)$ blocks, all that the adversary alone could have mined in $r$
rounds.}\label{fig.rates}
\end{figure}
